% ---------------------------------------------------------------------
%  10 Appendix
%  Included after \appendix in main.tex, so sections number as A, B, C, ...
%  \AppendixHeader (preamble.tex) prints the banner that opens the block.
% ---------------------------------------------------------------------
\AppendixHeader

\section{Notation}
\label{sec:appendix_notation}

\Todo{Notation table. Keep it generated from the macros in
\texttt{preamble.tex} so the symbols cannot drift.}

\begin{table}[h]
  \centering
  \begin{tabular}{ll}
    \toprule
    Symbol & Meaning \\
    \midrule
    $\numcams$, $\cam$ & number of cameras; camera index \\
    $\numframes$, $\timestep$ & sequence length; frame index \\
    $\homog{\cam}$, $\proj{\cam}$ & court homography; projection matrix \\
    $\ballTwoD{\cam,\timestep}$ & 2D ball detection in view $\cam$ \\
    $\ballThreeD{\timestep}$ & 3D ball position in court coordinates \\
    $\numcourtkp$ & number of court keypoints ($=14$) \\
    $\numjoints$ & number of body joints \\
    $(\rootrot{\timestep}, \rootpos{\timestep})$ & player root transform in court frame \\
    \bottomrule
  \end{tabular}
  \caption{Symbols used throughout the paper.}
  \label{tab:notation}
\end{table}

\section{Full training configuration}
\label{sec:appendix_config}
\Todo{Per-module hyperparameters, augmentation, schedules. Tables, not prose.}

\section{Reproduction}
\label{sec:appendix_repro}
\Todo{Checkpoint-to-result mapping and how to regenerate every number in
\cref{sec:results} from the released bundles.}

\section{Additional qualitative results}
\label{sec:appendix_qualitative}
\Todo{Extra reconstructions, including failure cases referenced in
\cref{subsec:failure_modes}.}
